%% Работа с русским языком и шрифтами (XeLaTeX)
\usepackage{fontspec}
\defaultfontfeatures{Ligatures=TeX}
\setmainfont{Times New Roman}[
  Path=C:/Windows/Fonts/,
  UprightFont=times.ttf,
  BoldFont=timesbd.ttf,
  ItalicFont=timesi.ttf,
  BoldItalicFont=timesbi.ttf
]
\setsansfont{Times New Roman}[
  Path=C:/Windows/Fonts/,
  UprightFont=times.ttf,
  BoldFont=timesbd.ttf,
  ItalicFont=timesi.ttf,
  BoldItalicFont=timesbi.ttf
]
\setmonofont{Courier New}[
  Path=C:/Windows/Fonts/,
  UprightFont=cour.ttf,
  BoldFont=courbd.ttf,
  ItalicFont=couri.ttf,
  BoldItalicFont=courbi.ttf,
  Scale=MatchLowercase
]
\usepackage[russian]{babel}
\usepackage{microtype}
\usepackage{xurl}
\usepackage{fvextra}

% Перенос длинных латинских идентификаторов в русском тексте (XeLaTeX)
\XeTeXlinebreaklocale "en"
\XeTeXlinebreakskip = 0pt plus 3pt minus 1pt

\RecustomVerbatimEnvironment{verbatim}{Verbatim}{
  breaklines=true,
  breakanywhere=true,
  breaksymbol={},
  fontsize=\small
}

%% Пакеты для работы с математикой
\usepackage{amsmath,amsfonts,amssymb,amsthm,mathtools}
\usepackage{icomma}

%% Нумерация формул (опционально)
%\mathtoolsset{showonlyrefs=true}
%\usepackage{leqno}

%% Шрифты
\usepackage{euscript}
\usepackage{mathrsfs}

%% Некоторые полезные макросы для дебага
\makeatletter
\newcommand\thefontsize{The current font size is: \f@size pt}
\makeatother

%% Настройка размеров шрифтов
\makeatletter
\setlength{\headheight}{28pt}
\renewcommand\Huge{\@setfontsize\Huge{14pt}{10}}
\renewcommand\huge{\@setfontsize\huge{14pt}{10}}
\renewcommand\Large{\@setfontsize\Large{14pt}{10}}
\renewcommand\large{\@setfontsize\large{12pt}{10}}
\makeatother

%% Поля (геометрия страницы)
\usepackage[left=3cm,right=1.5cm,top=2cm,bottom=2cm,bindingoffset=0cm]{geometry}

%% Русские списки
\usepackage{enumitem}
\makeatletter
\AddEnumerateCounter{\asbuk}{\russian@alph}{щ}
\makeatother

%% Работа с картинками
\usepackage{caption}
\captionsetup{labelsep=endash}
\captionsetup[figure]{justification=centering}
\captionsetup[table]{justification=raggedright, singlelinecheck=false}
\addto\captionsrussian{\renewcommand{\figurename}{Рисунок}}
\usepackage{graphicx}
\graphicspath{{images/}{images2/}}
\setlength\fboxsep{3pt}
\setlength\fboxrule{1pt}
\usepackage{wrapfig}

%% Работа с таблицами
\usepackage{array,tabularx,tabulary,booktabs}
\usepackage{longtable}
\usepackage{multirow}
\usepackage{makecell}
\usepackage{hhline}

%% Красная строка
\setlength{\parindent}{1.25cm}
\emergencystretch=4em
\tolerance=4000
\pretolerance=1000
\AtBeginDocument{\hbadness=10000}

%% Интервалы (ГОСТ: полуторный межстрочный интервал)
\usepackage{setspace}
\onehalfspacing

%% Компактные вертикальные отступы вокруг плавающих объектов, подписей и списков
\setlength{\textfloatsep}{8pt plus 2pt minus 2pt}
\setlength{\floatsep}{8pt plus 2pt minus 2pt}
\setlength{\intextsep}{8pt plus 2pt minus 2pt}
\captionsetup{skip=4pt}
\setlist{topsep=3pt, itemsep=2pt, parsep=0pt}

%% Заголовки
\usepackage{titlesec}
\titleformat{\section}[block]
  {\normalfont\bfseries\fontsize{14pt}{16pt}\selectfont\centering}
  {\thesection}{1em}{}
\titleformat{\subsection}[block]
  {\normalfont\bfseries\fontsize{14pt}{16pt}\selectfont\centering}
  {\thesubsection}{1em}{}
\titlespacing*{\section}{0pt}{1.6ex plus .4ex minus .2ex}{1.0ex plus .2ex}
\titlespacing*{\subsection}{0pt}{1.3ex plus .3ex minus .2ex}{0.8ex plus .2ex}

%% TikZ
\usepackage{tikz}
\usetikzlibrary{graphs,graphs.standard}
\usetikzlibrary{positioning,arrows.meta,shapes.geometric,shapes.misc,fit,backgrounds,calc}

%% Цветовая палитра для схем
\usepackage{xcolor}
\definecolor{schemeBlue}{HTML}{2C5F8A}
\definecolor{schemeBlueLt}{HTML}{D6E6F2}
\definecolor{schemeGreen}{HTML}{2E7D5B}
\definecolor{schemeGreenLt}{HTML}{D5EFE2}
\definecolor{schemeOrange}{HTML}{C2611F}
\definecolor{schemeOrangeLt}{HTML}{FBE3D0}
\definecolor{schemeGray}{HTML}{4A4A4A}
\definecolor{schemeGrayLt}{HTML}{ECECEC}
\definecolor{schemeRed}{HTML}{B23A3A}

%% Верхний и нижний колонтитулы
\usepackage{fancyhdr}
\pagestyle{fancy}
\fancyhf{}
\fancyfoot[C]{\thepage}
\renewcommand{\headrulewidth}{0pt}
\renewcommand{\footrulewidth}{0pt}
\fancypagestyle{plain}{
  \fancyhf{}
  \fancyfoot[C]{\thepage}
  \renewcommand{\headrulewidth}{0pt}
  \renewcommand{\footrulewidth}{0pt}
}

%% Перенос знаков в формулах (по Львовскому)
\newcommand*{\hm}[1]{#1\nobreak\discretionary{}{\hbox{$\mathsurround=0pt #1$}}{}}

%% Дополнительно
\usepackage{float}
\usepackage{gensymb}
\usepackage{listings}
\let\monospacetext\texttt
% \texttt переопределён выше (переносы длинных путей и идентификаторов)
% Переменные текстом (не math-mode): корректный слой PDF для Антиплагиата
\newcommand{\kreq}{K\_req}
\newcommand{\kmin}{K\_min}
\newcommand{\aptimes}[1]{×#1}
\lstset{
basicstyle=\small\ttfamily,
columns=flexible,
breaklines=true
}

% Hyperref (для ссылок внутри pdf; hidelinks — меньше шума в текстовом слое)
\usepackage[hidelinks,unicode]{hyperref}
\urlstyle{same}
\def\UrlFont{\normalfont}
\def\UrlBreaks{\do\/\do-\do\_\do\.\do\ \do\:\do\;\do\=\do\~\do\(\do\)\do\[\do\]}

% Идентификаторы и пути — обычный текст (переносы: XeTeX linebreak + emergencystretch)
\DeclareRobustCommand{\texttt}[1]{%
  \ifmmode
    \text{\normalfont #1}%
  \else
    {\normalfont #1}%
  \fi
}

% Термин — определение (единый формат, без скачков description)
\newcommand{\defitem}[2]{%
  \par\smallskip\noindent\textbf{#1}~---~#2%
}

% Отступ перед первым абзацем в каждом разделе
\usepackage{indentfirst}

% Глубина нумерации и оглавления
\setcounter{secnumdepth}{2}
\setcounter{tocdepth}{2}
